% Figure from MATH 493 notes, section 61. Compile with the notes preamble.
\begin{tikzpicture}[>=Stealth]
\draw[black!55, thin] (0,0) circle (2.0);
\draw[blue!70!black, very thick] (-2.349,-0.855) -- (2.349,0.855) node[pos=0.98, above, font=\scriptsize] {$L$};
\fill[blue!70!black] (1.879,0.684) circle (1.5pt);
\draw[orange!85!black, thick] (1.915,-1.607) -- (-1.915,1.607) node[pos=0.98, above, font=\scriptsize] {$gL$};
\fill[orange!85!black] (-1.532,1.286) circle (1.5pt);
\draw[red!70!black, thick] (0.434,2.462) -- (-0.434,-2.462) node[pos=0.98, right, xshift=4pt, font=\scriptsize] {$g^2L$};
\fill[red!70!black] (-0.347,-1.970) circle (1.5pt);
\draw[->, black!70] (20.0:1.2) arc (20.0:140.0:1.2) node[midway, above left, font=\scriptsize] {$120^\circ$};
\draw[->, black!70] (140.0:0.8) arc (140.0:260.0:0.8);
\draw[black!55, densely dotted] (-2.5,0) -- (2.5,0) node[right, font=\scriptsize] {$\mathbb{R}b_1$};
\draw[black!55, densely dotted] (0,-2.5) -- (0,2.5) node[above, font=\scriptsize] {$\mathbb{R}b_2$};
\node[font=\scriptsize, align=left, anchor=west] at (3.0,1.0) {$g = (1\,2\,3)$ turns $U$ by $120^\circ$, so it moves\\a line at angle $\alpha$ to the line at angle $\alpha + 120^\circ$;\\a line is unchanged only by turns of $0^\circ$ or $180^\circ$:\\\textcolor{blue!70!black}{$L$}, \textcolor{orange!85!black}{$gL$}, \textcolor{red!70!black}{$g^2L$} are three different lines.\\Dotted: the two lines that $(1\,2)$ alone preserves;\\\quad $(1\,2\,3)$ moves both.};
\end{tikzpicture}
